\subsubsection*{\texttt{10-baseline-no-rag}}
\label{app:samples:10_baseline_no_rag}

\paragraph{Q3: How long does the Haste spell last?}
\noindent \textbf{Corr.:}~2.00/5\quad \textbf{Faith.:}~—\quad \textbf{Relev.:}~—\quad \textbf{Ctx.Rel.:}~—\\

\noindent\textbf{Antwort:} The duration of the Haste spell depends on the level of the spell being cast. Here are the durations for different levels:  * Haste (1st level): 1 minute * Haste (2nd level): 1 minute * Haste (3rd level): 10 minutes * Haste (4th level): 1 hour * Haste (5th level): 8 hours * Haste (6th level): 24 hou\ldots

\vspace{4pt}
\hrule
\vspace{6pt}

\paragraph{Q5: Can a grappled creature cast spells?}
\noindent \textbf{Corr.:}~2.00/5\quad \textbf{Faith.:}~—\quad \textbf{Relev.:}~—\quad \textbf{Ctx.Rel.:}~—\\

\noindent\textbf{Antwort:} A grappled creature can still take actions, including casting spells, but the casting of the spell may be hindered by the grappling effect. The specifics would depend on the details of the situation and the spell being cast.  According to the rules in the Monster Manual (page 155), a grappled creatu\ldots

\vspace{4pt}
\hrule
\vspace{6pt}

\paragraph{Q12: If a creature has exhaustion level 3 and is also restrained, what is its movement speed?}
\noindent \textbf{Corr.:}~2.00/5\quad \textbf{Faith.:}~—\quad \textbf{Relev.:}~—\quad \textbf{Ctx.Rel.:}~—\\

\noindent\textbf{Antwort:} A creature with exhaustion level 3 has its movement speed reduced to 0 squares (5 feet in DnD 5e) per turn. The restrained condition does not affect a creature's movement speed directly. Therefore, the creature's movement speed is 0 squares while it has exhaustion level 3.

\vspace{4pt}
\hrule
\vspace{6pt}

\paragraph{Q16: Can you use Action Surge to cast two leveled spells in one turn?}
\noindent \textbf{Corr.:}~1.00/5\quad \textbf{Faith.:}~—\quad \textbf{Relev.:}~—\quad \textbf{Ctx.Rel.:}~—\\

\noindent\textbf{Antwort:} No, using Action Surge does not allow you to cast two leveled spells as part of the same turn. Action Surge is a feature that gives you an additional action on your turn, which you can use to take an additional action, such as casting a spell, making an attack, or using an ability. However, each act\ldots

\vspace{4pt}
\hrule
\vspace{6pt}
